%==============================================================================
% PARTIDA 05: CAJAS DE PASO Y DERIVACION
% Codigo: OE.5.2.4.1
%==============================================================================
\subsection{OE.5.2.4.1 Caja de Pase Rectangular 4''x2'' FG}

\subsubsection{Especificaciones Tecnicas}

Caja rectangular de fierro galvanizado (FG) de 4''x2'' para paso y derivacion de canalizaciones. Permite cambios de direccion de tuberias y acceso a conexiones de conductores para mantenimiento.

\begin{itemize}[nosep]
  \item \textbf{Material:} Fierro galvanizado por inmersion en caliente, espesor 1.2 mm.
  \item \textbf{Dimensiones:} 100 mm x 50 mm x 50 mm.
  \item \textbf{Entradas:} 4 knockouts de 3/4'' en caras laterales, 2 knockouts en base.
  \item \textbf{Tapa:} Metalica ciega, con tornillo de fijacion 1/4''-20 UNC.
  \item \textbf{Capacidad maxima de conductores:} Segun CNE-U Tabla 240.1, para caja 4''x2'':
  \begin{itemize}
    \item Conductores de 1.5 mm2: maximo 8 conductores.
    \item Conductores de 2.5 mm2: maximo 6 conductores.
  \end{itemize}
\end{itemize}

\subsubsection{Requisitos de Materiales}

\begin{longtable}{L{5.0cm} L{2.0cm} L{7.5cm}}
\toprule
\textbf{Material} & \textbf{Cant.} & \textbf{Especificacion tecnica} \\
\midrule
Caja rectangular FG 4''x2'' & 6 pza | FG galvanizado, con tapa ciega metalica \\
Conector "romex" 3/4'' & 12 pza | Metalico con contratuerca \\
Tornillo de anclaje 1/4'' x 2'' con tarugo & 12 pza | Para fijacion a muro \\
\bottomrule
\end{longtable}

\subsubsection{Metodo de Ejecucion}

\begin{enumerate}
  \item Ubicar la caja de paso en muro segun plano IE-04, en puntos donde la canalizacion requiera cambios de direccion.
  \item Empotrar en muro siguiendo el procedimiento de cajas rectangulares (OE.5.2.1.1.2).
  \item Conectar tuberias de entrada y salida mediante conectores "romex".
  \item Pasar los conductores a traves de la caja. No almacenar exceso de conductor dentro de la caja (maximo 15 cm de longitud adicional).
  \item Colocar la tapa ciega metalica y fijar con tornillo.
  \item Identificar los circuitos que pasan por la caja con etiqueta adhesiva en la parte interior de la tapa.
\end{enumerate}

\subsubsection{Control de Calidad}

\begin{itemize}
  \item Verificar que la caja quede al ras con el muro terminado (tolerancia: ±1 mm).
  \item Comprobar que las tuberias ingresen correctamente mediante conectores (no directamente).
  \item Verificar que la tapa cierre correctamente.
  \item Inspeccionar que los conductores tengan libertad de movimiento dentro de la caja.
\end{itemize}

\subsubsection{Metodo de Medicion}

Se medira por pieza (pza) de caja de paso rectangular completamente instalada.

\subsubsection{Unidad de Medida}

pza (pieza).

\subsubsection{Forma de Pago}

Pago por pieza instalada con tapa colocada y aprobada por supervision.

%==============================================================================
% PARTIDA 06: TABLEROS DE DISTRIBUCION
% Codigos: OE.5.2.6.1, OE.5.2.6.2, OE.5.2.6.3
%==============================================================================
\subsection{OE.5.2.6.1 Tablero General TG-01 (12 polos, metalico empotrable)}

\subsubsection{Especificaciones Tecnicas}

Tablero general de distribucion tipo empotrable para el primer piso. Debe cumplir con NTP 370.053, NTP-IEC 60439-1 (Conjuntos de aparamenta de baja tension) y CNE-U Articulo 380.1.

\begin{itemize}[nosep]
  \item \textbf{Tipo:} Empotrable en muro, con puerta ciega y cerradura.
  \item \textbf{Material:} Lamina de acero laminado en frio (LAC) de 1.5 mm de espesor, con pintura electrostática horneable color blanco (RAL 9010).
  \item \textbf{Capacidad:} 12 polos (12 modulos DIN de 17.5 mm cada uno = 210 mm utiles).
  \item \textbf{Riel DIN:} Riel simetrico de 35 mm (tipo Omega) segun NTP-IEC 60715.
  \item \textbf{Barras incluidas:}
  \begin{itemize}
    \item Barra de neutro: Aislada sobre soportes plasticos, con 12 bornes para conductores hasta 16 mm2.
    \item Barra de tierra: Conectada al gabinete, con 12 bornes para conductores hasta 16 mm2.
  \end{itemize}
  \item \textbf{Grado de proteccion:} IP-30 (interior), IK-07 (resistencia al impacto).
  \item \textbf{Acceso:} Puerta con bisagras y cerradura tipo manija con llave.
  \item \textbf{Dimensiones referenciales:} 300 mm (ancho) x 400 mm (alto) x 120 mm (fondo).
  \item \textbf{Capacidad de corriente:} 100 A maximo nominal del gabinete.
\end{itemize}

\paragraph{Circuitos alojados en TG-01:}
\begin{longtable}{L{2.5cm} L{5.0cm} L{3.0cm} L{3.5cm}}
\toprule
\textbf{Circuito} & \textbf{Proteccion} & \textbf{Conductor} & \textbf{Uso} \\
\midrule
General & Interruptor general 2P-32A & 10 mm2 & Alimentacion general \\
ID & Interruptor diferencial 2P-40A / 30mA & 10 mm2 & Proteccion de personas \\
C1 & Termomagnetico 2P-10A & 1.5 mm2 & Alumbrado primer piso \\
C2 & Termomagnetico 2P-16A + ID 25A & 2.5 mm2 & Tomacorrientes primer piso \\
\bottomrule
\end{longtable}

\subsubsection{Requisitos de Materiales}

\begin{longtable}{L{5.0cm} L{2.0cm} L{7.5cm}}
\toprule
\textbf{Material} & \textbf{Cant.} & \textbf{Especificacion tecnica} \\
\midrule
Tablero metalico 12 polos empotrable & 1 u | Acero 1.5 mm, pintura electrostatica, puerta con cerradura \\
Riel DIN simetrico 35 mm x 210 mm & 1 pza | NTP-IEC 60715, tipo Omega \\
Barra de neutro aislada 12 bornes & 1 pza | Con soportes aislantes, bornes para 16 mm2 \\
Barra de tierra 12 bornes & 1 pza | Conectada al gabinete, bornes para 16 mm2 \\
Interruptor termomagnetico 2P-32A curva C & 1 pza | Interruptor general \\
Interruptor diferencial 2P-40A / 30mA tipo A & 1 pza | Proteccion general de personas \\
Interruptor termomagnetico 2P-10A curva C & 1 pza | Proteccion C1 \\
Interruptor termomagnetico 2P-16A curva C & 1 pza | Proteccion C2 \\
Interruptor diferencial 2P-25A / 30mA tipo A & 1 pza | Proteccion especifica C2 \\
Separadores de conductores plasticos & 4 pza | Para organizacion interior del tablero \\
Etiquetas rotuladoras de circuitos & 1 juego | Adhesivas, para identificar cada circuito \\
Canaleta ranurada 25x25 mm & 1 m | Para organizacion de cableado interno \\
\bottomrule
\end{longtable}

\subsubsection{Metodo de Ejecucion}

\begin{enumerate}
  \item \textbf{Replanteo:} Ubicar TG-01 en primer piso, cerca del ingreso, en zona de facil acceso (CNE-U Articulo 380.2). Altura: 1.60 m sobre NPT al centro del tablero.
  \item \textbf{Apertura de muro:} Realizar la excavacion para empotrar el tablero (dimensiones: 320 mm x 420 mm x 130 mm).
  \item \textbf{Fijacion:} Colocar el tablero en el hueco y fijar con pernos de expansion de 1/4'' x 2'' en las 4 esquinas internas. Nivelar con nivel de burbuja (tolerancia: ±2 mm en ambos ejes).
  \item \textbf{Instalacion de componentes internos:}
  \begin{enumerate}
    \item Montar el riel DIN en los soportes del gabinete.
    \item Montar las barras de neutro y tierra en sus soportes aislantes.
    \item Insertar los interruptores en el riel DIN (encajar a presion en la posicion asignada).
  \end{enumerate}
  \item \textbf{Cableado:}
  \begin{enumerate}
    \item Conectar el alimentador general (10 mm2) a la entrada del interruptor general (borne superior).
    \item Conectar la salida del interruptor general a la entrada del interruptor diferencial general.
    \item Conectar la salida del ID general a la barra de distribucion de fase.
    \item Conectar cada circuito derivado a su respectivo interruptor termomagnetico.
    \item Conectar los neutros a la barra de neutro.
    \item Conectar los conductores de proteccion a la barra de tierra.
  \end{enumerate}
  \item \textbf{Rotulado:} Colocar etiquetas adhesivas en cada interruptor indicando el circuito y su descripcion (ej: "C1 -- Alumbrado 1er Piso").
  \item \textbf{Organizacion:} Fijar los conductores en el interior del tablero utilizando canaletas ranuradas y separadores. Mantener separacion entre conductores de fase, neutro y tierra.
\end{enumerate}

\subsubsection{Control de Calidad}

\begin{longtable}{L{4.5cm} L{4.5cm} L{4.5cm}}
\toprule
\textbf{Parametro} & \textbf{Metodo de verificacion} & \textbf{Criterio de aceptacion} \\
\midrule
Nivelacion & Nivel de burbuja sobre puerta del tablero & ±2 mm en ambos ejes \\
Fijacion & Prueba de traccion manual (10 kg) & Sin desplazamiento \\
Separacion barra N/T & Inspeccion visual & Barra de neutro aislada, barra de tierra conectada a gabinete \\
Torque de conexiones & Llave de torque dinamometrica & Bornes de interruptores: 3.5 Nm; bornes de barras: 5 Nm \\
Rotulado & Inspeccion visual & Todos los circuitos identificados claramente \\
Continuidad de tierra & Multimetro entre barra de tierra y gabinete & < 0.5 $\Omega$ \\
\bottomrule
\end{longtable}

\subsubsection{Metodo de Medicion}

Se medira por unidad (u) de tablero completamente instalado, incluyendo gabinete, riel DIN, barras, interruptores, cableado interno, rotulado y pruebas.

\subsubsection{Unidad de Medida}

u (unidad).

\subsubsection{Forma de Pago}

Pago por unidad instalada, cableada, rotulada, probada y aprobada por supervision.

%------------------------------------------------------------------------------
\subsection{OE.5.2.6.2 Tablero de Distribucion TD-01 (8 polos, segundo piso)}

\subsubsection{Especificaciones Tecnicas}

Tablero de distribucion secundario empotrable de 8 polos para el segundo piso.

\begin{itemize}[nosep]
  \item Capacidad: 8 polos (8 modulos DIN = 140 mm utiles).
  \item Dimensiones referenciales: 250 mm x 300 mm x 100 mm.
  \item Circuitos alojados: C3 (alumbrado 2do piso, 2P-10A) y C4 (tomacorrientes 2do piso, 2P-16A + ID 25A).
  \item Alimentacion: Desde TG-01 con conductores de 6 mm2 (fase + neutro + tierra).
\end{itemize}

Metodo de ejecucion, control de calidad y forma de pago identicos a OE.5.2.6.1.

\subsubsection{Unidad de Medida}

u (unidad).

%------------------------------------------------------------------------------
\subsection{OE.5.2.6.3 Tablero de Distribucion TD-02 (8 polos, tercer piso)}

Identico a TD-01 (OE.5.2.6.2) pero ubicado en el tercer piso. Circuitos alojados: C5 (alumbrado 3er piso) y C6 (tomacorrientes 3er piso).

\subsubsection{Unidad de Medida}

u (unidad).

%==============================================================================
% PARTIDA 07: DISPOSITIVOS DE MANIOBRA Y PROTECCION
% Codigos: OE.5.2.8.1 al OE.5.2.8.7
%==============================================================================
\subsection{OE.5.2.8.1 Interruptor Simple Unipolar (10 A)}

\subsubsection{Especificaciones Tecnicas}

Interruptor simple unipolar (S1) de 10 A -- 250 V, para control de iluminacion desde una sola ubicacion.

\begin{itemize}[nosep]
  \item \textbf{Corriente nominal:} 10 A.
  \item \textbf{Tension nominal:} 250 V.
  \item \textbf{Material de contactos:} Aleacion de plata-cadmio (Ag-CdO) para resistencia a la formacion de arco.
  \item \textbf{Material del cuerpo:} Polimero termoplastico autoextinguible (V-0 segun UL 94).
  \item \textbf{Vida util electrica:} Minimo 20000 ciclos de operacion.
  \item \textbf{Placa frontal:} 1 modulo, color blanco, ABS ignifugo.
  \item \textbf{Normas:} NTP-IEC 60669-1 (Interruptores para instalaciones electricas fijas domesticas).
\end{itemize}

\paragraph{Conexion electrica:}
\begin{itemize}[nosep]
  \item Borne comun (L): Conexion de fase de alimentacion.
  \item Borne de retorno ($\uparrow$): Conexion a luminaria.
  \item No conectar neutro al interruptor.
\end{itemize}

\subsubsection{Metodo de Ejecucion}

\begin{enumerate}
  \item Identificar el conductor de fase (rojo/marron) y el conductor de retorno (debe ser del mismo color de fase, identificado con cinta).
  \item Pelar 1.2 cm de aislamiento en cada conductor.
  \item Conectar fase al borne comun (L) del interruptor.
  \item Conectar retorno al borne de salida ($\uparrow$).
  \item Fijar el mecanismo a la caja rectangular mediante los tornillos de las orejas.
  \item Colocar la placa frontal decorativa.
\end{enumerate}

\subsubsection{Control de Calidad}

\begin{itemize}
  \item Verificar que el interruptor controle la fase (no el neutro).
  \item Probar funcionamiento (debe encender y apagar la luminaria correctamente).
  \item Verificar que la placa frontal quede al ras con el muro terminado.
\end{itemize}

\subsubsection{Metodo de Medicion}

Se medira por pieza (pza) de interruptor simple instalado y probado.

\subsubsection{Unidad de Medida}

pza (pieza).

%------------------------------------------------------------------------------
\subsection{OE.5.2.8.2 Interruptor de Conmutacion 3 Vias (S3)}

\subsubsection{Especificaciones Tecnicas}

Interruptor de conmutacion de 3 vias (S3) para control de luminaria desde dos ubicaciones (escaleras, pasadizos).

\begin{itemize}[nosep]
  \item Corriente nominal: 10 A.
  \item Tension nominal: 250 V.
  \item Bornes: 3 bornes (L/comun, 1 viajero, 2 viajero).
  \item Norma: NTP-IEC 60669-1.
\end{itemize}

\paragraph{Cableado tipico para escaleras:}
\begin{enumerate}
  \item Primer S3 (parte inferior de escalera): Recibe fase en borne comun.
  \item Segundo S3 (parte superior de escalera): Borne comun conectado al retorno de la luminaria.
  \item Conductores viajeros conectan bornes 1 y 2 de ambos interruptores.
\end{enumerate}

\subsubsection{Metodo de Ejecucion}

\begin{enumerate}
  \item Identificar bornes: buscar el borne comun (generalmente marcado como "L", "C" o "COM").
  \item Conectar fase al borne comun del primer S3.
  \item Conectar dos conductores viajeros entre los bornes 1 y 2 de ambos interruptores.
  \item Conectar el borne comun del segundo S3 al conductor de retorno de la luminaria.
\end{enumerate}

\subsubsection{Control de Calidad}

\begin{itemize}
  \item Verificar que la luminaria se encienda/apague desde ambos interruptores en cualquier combinacion.
  \item Probar todas las posiciones de los interruptores.
\end{itemize}

\subsubsection{Unidad de Medida}

pza (pieza).

%------------------------------------------------------------------------------
\subsection{OE.5.2.8.3 al OE.5.2.8.7 Interruptores Termomagneticos y Diferencial}

\subsubsection{Especificaciones Tecnicas}

\begin{longtable}{L{2.5cm} L{3.0cm} L{2.5cm} L{2.0cm} L{4.5cm}}
\toprule
\textbf{Codigo} & \textbf{Tipo} & \textbf{Especificacion} & \textbf{Cant.} & \textbf{Uso} \\
\midrule
OE.5.2.8.3 & Termomagnetico 2P-10A & 10 A, curva C, 10 kA, 230/400 V & 3 & Proteccion alumbrado C1, C3, C5 \\
OE.5.2.8.4 & Termomagnetico 2P-16A & 16 A, curva C, 10 kA, 230/400 V & 3 & Proteccion tomacorrientes C2, C4, C6 \\
OE.5.2.8.5 & Termomagnetico 2P-20A & 20 A, curva C, 10 kA, 230/400 V & 1 & Proteccion cocina / cargas especiales \\
OE.5.2.8.6 & Termomagnetico general 2P-40A & 40 A, curva C, 10 kA, 230/400 V & 1 & Interruptor general TG-01 \\
OE.5.2.8.7 & Diferencial 2P-25A-30mA & 25 A, 30 mA, tipo A, 230 V & 1 & Proteccion de personas (ID general) \\
\bottomrule
\end{longtable}

\paragraph{Caracteristicas de los interruptores termomagneticos:}
\begin{itemize}[nosep]
  \item Curva de disparo: Tipo C (disparo magnetico entre 5 y 10 x In), para cargas residenciales.
  \item Poder de corte: 10 kA segun NTP-IEC 60898-1.
  \item Tension nominal: 230/400 V.
  \item Conexion: Bornes tipo jaula para conductores hasta 25 mm2 (general) y 16 mm2 (derivados).
  \item Indicacion de posicion: Rojo/Verde (ON/OFF).
  \item Norma: NTP-IEC 60898-1 (Interruptores automaticos para proteccion contra sobrecorrientes).
\end{itemize}

\paragraph{Caracteristicas del interruptor diferencial:}
\begin{itemize}[nosep]
  \item Corriente nominal ($I_n$): 25 A.
  \item Corriente diferencial residual ($I_{\Delta n}$): 30 mA.
  \item Tipo: A (protege corrientes alternas senoidales y pulsantes).
  \item Tiempo de disparo: < 40 ms a 1 x $I_{\Delta n}$.
  \item Poder de corte: 10 kA.
  \item Norma: NTP-IEC 61008-1 (Interruptores automaticos para operar por corriente diferencial residual).
\end{itemize}

\subsubsection{Metodo de Ejecucion}

\begin{enumerate}
  \item Montar cada interruptor en el riel DIN del tablero, encajando a presion.
  \item Para interruptores termomagneticos:
  \begin{enumerate}
    \item Conectar fase de entrada al borne superior.
    \item Conectar fase de salida (hacia la carga) al borne inferior.
    \item No conectar neutro al termomagnetico (es unipolar por fase).
  \end{enumerate}
  \item Para interruptor diferencial (2 polos):
  \begin{enumerate}
    \item Conectar fase y neutro de entrada en bornes superiores (marcados como LINE o 1-2).
    \item Conectar fase y neutro de salida en bornes inferiores (marcados como LOAD o 3-4).
    \item Respetar la polaridad: fase en el borne izquierdo, neutro en el derecho.
  \end{enumerate}
  \item Torque de apriete: 3.5 Nm para conductores de 2.5-6 mm2; 5 Nm para conductores de 10 mm2.
\end{enumerate}

\subsubsection{Control de Calidad}

\begin{longtable}{L{4.5cm} L{4.5cm} L{4.5cm}}
\toprule
\textbf{Parametro} & \textbf{Metodo de verificacion} & \textbf{Criterio de aceptacion} \\
\midrule
Corriente nominal & Inspeccion visual de la marcacion en el interruptor & Coincide con la especificacion \\
Curva de disparo & Inspeccion visual de la marcacion & Tipo C \\
Poder de corte & Inspeccion visual & Minimo 6 kA (segun NTP-IEC 60898-1) \\
Funcionamiento termomagnetico & Prueba de accion manual (ON/OFF) & Mecanismo firme, sin juego \\
Funcionamiento diferencial & Prueba con boton TEST & Disparo inmediato al presionar TEST \\
Tiempo de disparo diferencial & RCD tester & < 40 ms a 30 mA \\
Torque de apriete & Llave de torque dinamometrica & 3.5 Nm (conductores 2.5-6 mm2) \\
\bottomrule
\end{longtable}

\subsubsection{Metodo de Medicion}

Se medira por pieza (pza) de interruptor instalado y probado.

\subsubsection{Unidad de Medida}

pza (pieza).

\subsubsection{Forma de Pago}

Pago por pieza instalada, probada (incluyendo prueba TEST en diferenciales) y aprobada por supervision.
